% Notation
% ~3 pages

\chapter*{표기법}
\addcontentsline{toc}{chapter}{표기법}
\markboth{표기법}{표기법}

이 책에서 사용하는 주요 기호와 규약을 정리한다. 각 기호는 본문에서 처음 등장할 때 다시 정의된다.


\section*{기본 표기}

\begin{tabular}{lp{11cm}}
\hline
\textbf{기호} & \textbf{의미} \\
\hline
$T$ & 시간 지표 집합 (전순서 집합) \\
$t, s \in T$ & 시간 지표 \\
$X_t$ & 시각 $t$에서의 관계적 지지 공간 (유한 집합) \\
$n = |X_t|$ & 지지 공간의 크기 \\
$x, y \in X_t$ & 지지 공간의 지점 \\
$u_t : X_t \to [0,1]$ & 부드러운 응집 장 (이론의 원초적 존재자) \\
$u_t(x)$ & 지점 $x$에서의 응집적 참여 강도 \\
$v_t = u_t - c\,\mathbf{1}$ & 평균 제거된 장 ($c = m/n$) \\
$\mathbf{1}$ & 모든 성분이 1인 벡터 \\
$[0,1]^{X_t}$ & $X_t$ 위의 $[0,1]$-값 함수들의 공간 \\
$\mathbf{u} = (u_t)_{t \in W}$ & 시간 창 $W \subseteq T$에 걸친 장의 모음 \\
\hline
\end{tabular}


\section*{형식 우주}

SCC 이론의 형식 우주는 다음 튜플이다:
$$
\mathfrak{C}^{\mathrm{soft}} = \Big( T,\ \{X_t\}_{t \in T},\ \{u_t\}_{t \in T},\ \{\mathrm{Cl}_t\}_{t \in T},\ \{\mathbf{N}_t, \mathbf{D}_t\}_{t \in T},\ \{\mathbf{M}_{t \to s}\}_{t,s \in T} \Big)
$$

\begin{tabular}{lp{11cm}}
\hline
\textbf{기호} & \textbf{의미} \\
\hline
$\mathrm{Cl}_t : [0,1]^{X_t} \to [0,1]^{X_t}$ & 부드러운 닫힘 연산자 (관계적 자기-완성) \\
$\mathbf{N}_t : X_t \times X_t \to [0,\infty)$ & 부드러운 인접 핵 (국소적 관계 구조) \\
$\mathbf{D}_t : X_t \times [0,1]^{X_t} \to [0,1]$ & 부드러운 구별 연산자 (내외부 비대칭) \\
$\mathbf{C}_t : X_t \times X_t \to [0,\infty)$ & 공동-소속 연산자 (파생 진단, 형식 우주 외부) \\
$\mathbf{M}_{t \to s} : X_t \times X_s \to [0,1]$ & 시간적 전달 핵 (응집 조직의 시간적 계승) \\
\hline
\end{tabular}


\section*{부피 제약과 에너지}

\begin{tabular}{lp{11cm}}
\hline
\textbf{기호} & \textbf{의미} \\
\hline
$m = \sum_{x} u_t(x)$ & 총 응집 부피 \\
$c = m/n$ & 부피 분율 (필드 평균값) \\
$\Sigma_m$ & 부피 제약 단체: $\{u \in [0,1]^n : \sum u_i = m\}$ \\[6pt]
$\mathcal{E}(\mathbf{u})$ & 총 에너지 범함수 \\
$\mathcal{E}_{\mathrm{cl}}$ & 닫힘 에너지: $\sum_x (u(x) - \mathrm{Cl}(u)(x))^2$ \\
$\mathcal{E}_{\mathrm{sep}}$ & 분리 에너지: $\sum_x u(x)(1 - \mathbf{D}(x; 1-u))$ \\
$\mathcal{E}_{\mathrm{bd}}$ & 경계/형태학 에너지: $\alpha \sum_{x,y} \mathbf{N}(x,y)(u(x)-u(y))^2 + \beta \sum_x W(u(x))$ \\
$\mathcal{E}_{\mathrm{tr}}$ & 전달 에너지: $\sum_{x,y} \mathbf{M}(x,y)\, \omega(u_t(x), u_s(y))(u_s(y) - u_t(x))^2$ \\[6pt]
$\lambda_{\mathrm{cl}}, \lambda_{\mathrm{sep}}, \lambda_{\mathrm{bd}}, \lambda_{\mathrm{tr}}$ & 에너지 가중 계수 \\
$\alpha, \beta$ & 형태학 매개변수 (매끄러움, 이중 우물) \\
$W(u) = u^2(1-u)^2$ & 이중 우물 포텐셜 \\
$W'(u) = 2u(1-u)(1-2u)$ & 이중 우물 기울기 (인자 2 포함) \\
\hline
\end{tabular}


\section*{진단 벡터}

\begin{tabular}{lp{11cm}}
\hline
\textbf{기호} & \textbf{의미} \\
\hline
$\mathbf{d} = (\mathsf{Bind}, \mathsf{Sep}, \mathsf{Inside}, \mathsf{Persist})$ & Proto-cohesion 진단 벡터 $\in [0,1]^4$ \\[4pt]
$\mathsf{Bind}_t(u_t) = 1 - \|u - \mathrm{Cl}(u)\|_2 / \sqrt{n}$ & 결합: 관계적 자기-지지의 정도 \\[4pt]
$\mathsf{Sep}_t(u_t) = \sum u \cdot \mathbf{D}(x;1\!-\!u) / \sum u$ & 분리: 외부와의 구조적 대비 \\[4pt]
$\mathsf{Inside}_t(u_t) = \mathcal{Q}_{\mathrm{morph}}(u_t)$ & 내부성: 형태학적 분절의 질 \\[4pt]
$\mathsf{Persist}(u_{\mathrm{prev}}, u_{\mathrm{curr}})$ & 지속성: 시간적 구조 계승 \\[6pt]
$\ell_{\max}$ & $H_0$ 지속성 다이어그램의 가장 긴 막대 길이 \\
$\mathrm{Artic}(u) = 1 - \ell_{\mathrm{second}}/\ell_{\max}$ & 분절 비율 \\
$\mathcal{Q}_{\mathrm{morph}} = \frac{\ell_{\max} - c}{1 - c} \cdot \mathrm{Artic}$ & 형태학적 질 측도 \\
\hline
\end{tabular}


\section*{그래프 구조}

\begin{tabular}{lp{11cm}}
\hline
\textbf{기호} & \textbf{의미} \\
\hline
$A$ & 인접 행렬: $A_{ij} = \mathbf{N}_t(x_i, x_j)$ \\
$D$ & 차수 행렬: $D_{ii} = \sum_j A_{ij}$ \\
$L = D - A$ & 그래프 라플라시안 (비정규화) \\
$P = D^{-1}A$ & 행 정규화 인접 행렬 (확률 전이 행렬) \\
$W_{\mathrm{sym}}$ & 응집 가중 대칭 행렬: $W_{ij} = A_{ij} \cdot u_i \cdot u_j$ \\
$\lambda_1 \leq \lambda_2 \leq \cdots$ & 라플라시안 고유값 ($\lambda_1 = 0$) \\
$\lambda_2$ & Fiedler 고유값 (대수적 연결도) \\
$a_{\mathrm{cl}} = \max_i P_{ii}$ & 닫힘 매개변수 상한 (공리: $a_{\mathrm{cl}} < 4$) \\
\hline
\end{tabular}


\section*{다중 형성}

\begin{tabular}{lp{11cm}}
\hline
\textbf{기호} & \textbf{의미} \\
\hline
$K$ & 형성의 수 (K-필드 아키텍처) \\
$u^{(k)}$ & $k$번째 형성의 응집 장 ($k = 1, \ldots, K$) \\
$\Lambda_{\mathrm{coupling}}$ & 결합 매개변수: $\lambda_{\mathrm{rep}} \cdot \omega_{jk} / \min(\mu_j, \mu_k)$ \\
$\lambda_{\mathrm{rep}}$ & 형성 간 반발(repulsion) 가중치 \\
$\omega_{jk}$ & 형성 $j$와 $k$ 사이의 공간적 중첩 \\
$\mu_k = \sum_x u^{(k)}(x)$ & $k$번째 형성의 부피 \\
$\mu_{\mathrm{floor}}$ & 부피 하한 정규화: $w_{\mathrm{cl}} \cdot 2(1 - a_{\mathrm{cl}}/4)^2$ \\
$d_{\min}^*$ & 최소 형성 간 거리 (안정성 임계) \\
\hline
\end{tabular}


\section*{파생 기하 개념}

\begin{tabular}{lp{11cm}}
\hline
\textbf{기호} & \textbf{의미} \\
\hline
$\mathrm{Core}_t(u_t) = \{x : u_t(x) \geq \theta_{\mathrm{core}}\}$ & 형성의 핵심부 \\
$\mathrm{Int}_t(u_t) = \{x : u_t(x) \geq \theta_{\mathrm{in}}\}$ & 형성의 내부 \\
$\mathrm{Bd}_t(u_t) = \{x : \theta_1 < u_t(x) < \theta_2\}$ & 경계 대역 \\
$g_t(x; u) = \sum_y \mathbf{N}_t(x,y)|u(x) - u(y)|$ & 기울기 지표 \\
\hline
\end{tabular}


\section*{합 규약}

이 책은 Canonical Specification v2.1의 합 규약을 따른다:

\begin{itemize}
    \item $\displaystyle\sum_{x,y \in X_t}$는 \textbf{순서쌍}(ordered pairs)을 범위로 한다. 대칭 핵에서 각 비방향 변은 두 번 세어진다. 이 규약은 위상 전이 정리(T8-Core)와 Hessian 분석에서 핵심적이다.
    \item $\displaystyle\sum_{\{x,y\}}$는 \textbf{비순서쌍}(unordered pairs)을 범위로 한다.
    \item 두 규약 사이의 관계: 대칭 함수 $f$에 대해 $\displaystyle\sum_{\{x,y\}} f(x,y) = \frac{1}{2}\sum_{x,y} f(x,y)$.
\end{itemize}

\textbf{이 규약은 핵심적이다.} 순서쌍/비순서쌍 혼동은 에너지 기울기, Hessian, 위상 전이 임계비에서 인자 2 오류를 발생시킨다. 의심스러운 경우 항상 순서쌍 기준으로 돌아가라.


\section*{노름과 거리}

\begin{tabular}{lp{11cm}}
\hline
\textbf{기호} & \textbf{의미} \\
\hline
$\|v\|_2 = \big(\sum_i v_i^2\big)^{1/2}$ & $\ell^2$ (유클리드) 노름 \\
$\|v\|_\infty = \max_i |v_i|$ & $\ell^\infty$ (최대) 노름 \\
$v^T L v$ & 그래프 Dirichlet 에너지 ($L$: 라플라시안) \\
\hline
\end{tabular}


\section*{자주 사용되는 약자}

\begin{tabular}{lp{11cm}}
\hline
\textbf{약자} & \textbf{의미} \\
\hline
SCC & Soft Cognitive Cohesion \\
Cat A/B/C & 정리 카테고리 (A: 완전 증명, B: 조건부, C: 추측) \\
FD & 유한차분 (finite difference)---기울기 검증에 사용 \\
PGD & 사영 기울기 하강 (projected gradient descent) \\
OT & 최적 전달 (optimal transport) \\
BB & Barzilai--Borwein 보폭 \\
$H_0$ & 0차 상동성 (연결 성분) \\
\hline
\end{tabular}


\section*{인덱싱 규약}

\begin{itemize}
    \item $t, s$: 시간 지표
    \item $x, y$: 공간 지표 (지지 공간의 지점)
    \item $i, j$: 지점의 정수 색인 ($x_i \in X_t$)
    \item $k$: 형성 색인 (다중 형성에서 $k = 1, \ldots, K$)
\end{itemize}
